% Generated draft: requires mathematical and source review.
\documentclass[11pt]{article}
\usepackage[margin=25mm]{geometry}
\usepackage{fontspec}
\setmainfont{Latin Modern Roman}
\usepackage{amsmath,amssymb,mathtools,unicode-math}
\setmathfont{Latin Modern Math}
\usepackage{enumitem,needspace,xurl,hyperref}
\setlength{\parindent}{0pt}
\setlength{\parskip}{5pt}
\setlength{\emergencystretch}{3em}
\newcommand{\N}{\mathbb{N}}
\newcommand{\Z}{\mathbb{Z}}
\newcommand{\Q}{\mathbb{Q}}
\newcommand{\R}{\mathbb{R}}
\newcommand{\C}{\mathbb{C}}
\newcommand{\F}{\mathbb{F}}
\newcommand{\A}{\mathbb{A}}
\newcommand{\PP}{\mathbb{P}}
\newcommand{\E}{\mathbb{E}}
\newcommand{\eps}{\varepsilon}
\newcommand{\dd}{\,\mathrm{d}}
\begin{document}
\section*{UnsolvedMath ID: 800903}
\textbf{Source problem code:} catalogue entry
\subsection{Definitions and mathematical statement}
For notation, on a complex manifold \(X\) with a smooth real \((1,1)\)-form \(\omega\), an \(\omega\)-plurisubharmonic function is an upper semicontinuous, locally integrable function \(\varphi:X\to[-\infty,\infty)\) satisfying
\[
\omega+dd^c\varphi\geq 0
\]
in the sense of currents, where \(dd^c=\frac{i}{\pi}\partial\bar\partial\). For such a function, its Lelong number at \(x\), normalized so that \(\log\lVert z\rVert\) has Lelong number one at the origin, is
\[
\nu(\varphi,x)=\liminf_{\substack{z\to x\\z\ne x}}\frac{\varphi(z)}{\log\lVert z-x\rVert}
\]
in local coordinates centered at \(x\). Equivalently, one may first add a smooth local function making \(\varphi\) plurisubharmonic; this does not change the Lelong number.

The supplied excerpt does not state the ambient hypotheses on \(X\) and \(\omega\), define \(\mathcal E(X,\omega)\), or specify the meaning of finite energy. The notation above explains the local terminology; it does not supply those missing hypotheses. Retain \(\mathcal E(X,\omega)\) as the source-designated class and finite energy as the source-designated condition.

\textbf{Question 34, attributed in the source to Berman--Boucksom--Eyssidieux--Guedj--Zeriahi.} Under the ambient hypotheses and energy convention intended by the source, is the following assertion true?
\[
\forall\varphi\in\mathcal E(X,\omega),\qquad
\bigl[\varphi\text{ is }\omega\text{-plurisubharmonic and has finite energy}\bigr]
\Longrightarrow
\bigl[\forall x\in X,\ \nu(\varphi,x)=0\bigr].
\]
A fully specified mathematical formulation requires recovery of the omitted ambient hypotheses and energy definitions. In particular, the conclusion concerns every point of \(X\), with no exceptional locus stated in the supplied excerpt.
\subsection{Short English statement}
Question 34 asks whether every finite-energy \(\omega\)-plurisubharmonic function in the source's class \(\mathcal E(X,\omega)\) has zero Lelong number at every point of \(X\). The ambient assumptions and the energy convention are missing from the supplied excerpt, so this is a conditional statement-only draft requiring mathematical review.
\subsection{Sources}
\begin{itemize}
\item Supplied local UnsolvedMath registry record: permanent ID \(800903\), exact source problem code \texttt{catalogue entry}, title \emph{Open Problems in Pluripotential Theory --- Question 34}. The registry lists Dinew, Guedj, and Zeriahi as proposers, with proposed year 2015.
\item Dinew, Guedj, and Zeriahi, \emph{Open Problems in Pluripotential Theory} (2015), Question 34. Supplied source URL: \url{https://arxiv.org/abs/1511.00705}. The supplied background records access on 2026-10-16 and extraction by \texttt{arxiv\_latex}; these are supplied provenance details, not a claim of independent source access or verification.
\item The supplied question explicitly attributes the problem to Berman--Boucksom--Eyssidieux--Guedj--Zeriahi. No separate bibliographic reference for that attribution is supplied.
\item The local registry labels the problem \texttt{open}. This is a supplied status label, not an independently verified literature assessment.
\end{itemize}
\end{document}
